\chapter{2023年全国乙卷（理）}

\section{单选题}

\begin{problem}
设 \(z=\frac{2+\i}{2+2\i}\)，则 \(z=\)（\quad）

\choices
    {\(1-\i\)}
    {\(1+\i\)}
    {\(2-\i\)}
    {\(2+\i\)}
\end{problem}

\begin{problem}
设集合 \(U=\R\)，集合 \(M=\{x\mid x<1\}\)，\(N=\{x\mid -1<x<2\}\)，则 \(M\cup N^c\text{ 的补集}\) 为（\quad）

\choices
    {\(\complement_U(M\cup N)\)}
    {\(N\cup\complement_U M\)}
    {\(\complement_U(M\cap N)\)}
    {\(M\cup\complement_U N\)}
\end{problem}

\begin{problem}
如图，网格纸上绘制的一个零件的三视图，网格小正方形边长为 1，则该零件表面积为（\quad）

\begin{center}
\FigureLayoutDeclare{img/2023/national_paper_b_science/q3_fig1.png}{6.0cm}{0bp 29bp 22bp 10bp}
\bitmapfigure[max width=6.0cm]{img/2023/national_paper_b_science/q3_fig1.png}
\end{center}

\choices
    {\(24\)}
    {\(26\)}
    {\(28\)}
    {\(30\)}
\end{problem}

\begin{problem}
已知 \(f(x)=\frac{x e^x}{e^{ax}-1}\)，若 \(f(x)\) 为偶函数，则 \(a=\)（\quad）

\choices
    {\(-2\)}
    {\(-1\)}
    {\(1\)}
    {\(2\)}
\end{problem}

\begin{problem}
设 \(O\) 为平面坐标系原点，在区域 \(1\le x+y\le4\) 内随机取一点 A. 则直线 \(OA\) 的倾斜角不大于 \(\frac\pi4\) 的概率为（\quad）

\choices
    {\(\frac12\)}
    {\(\frac16\)}
    {\(\frac14\)}
    {\(\frac12\)}
\end{problem}

\begin{problem}
已知 \(f(x)=\sin(\omega x+\varphi)\) 在区间 \(\left(-\frac\pi6,\frac\pi3\right)\) 内图形关于 \(x=-\frac1\omega\) 对称，则 \(\omega=\)（\quad）

\choices
    {\(-\frac32\)}
    {\(-\frac12\)}
    {\(\frac12\)}
    {\(\frac32\)}
\end{problem}

\begin{problem}
甲乙两位同学从 6 种课外读物中各自选读 2 种，则这两人选读的课外读物中恰有 1 种相同的选法共有（\quad）

\choices
    {30}
    {60}
    {120}
    {240}
\end{problem}

\begin{problem}
已知圆锥 \(PO\) 的底面半径为 3，\(O\) 为底面圆心，\(PA\)，\(PB\) 为圆锥母线，\(\angle AOB=120^\circ\)，若 \(\triangle PAB\) 面积为 \( \frac{3\sqrt3}{4}\)，则该圆锥体积为（\quad）

\choices
    {\(\pi\)}
    {\(6\pi\)}
    {\(3\pi\)}
    {\(\frac{3\pi}{2}\)}
\end{problem}

\begin{problem}
已知 \(\triangle ABC\) 为等腰直角三角形，\(AB\) 为斜边，\(\triangle ABD\) 为等边三角形，若二面角 \(C-AB-D\) 为 \(150^\circ\)，则直线 \(CD\) 与平面 \(ABC\) 所成角的正切值为（\quad）

\choices
    {\(\frac15\)}
    {\(\frac{3}{5}\)}
    {\(\frac{5}{3}\)}
    {\(\frac{5}{\sqrt3}\)}
\end{problem}

\begin{problem}
已知等差数列 \(\{a_n\}\) 的公差为 \(\frac{2\pi}{3}\)，集合 \(S=\{\cos a_n\mid n\in\N^*\}=\{a,b\}\)，则 \(ab=\)（\quad）

\choices
    {\(-1\)}
    {\(-\frac12\)}
    {0}
    {\(\frac12\)}
\end{problem}

\begin{problem}
设 \(A\)，\(B\) 为双曲线 \(\frac{x^2}{9}-\frac{y^2}{1}=1\) 上两点，下列点中可为 \(\overline{AB}\) 中点的是（\quad）

\choices
    {\((1,1)\)}
    {\((-1,2)\)}
    {\((1,3)\)}
    {\((-1,-4)\)}
\end{problem}

\begin{problem}
已知圆 \(O\) 的半径为 1，直线 \(PA\) 与 \(O\) 相切于 \(A\)，\(\overline{PB}\cap\overline{PC}\) 与 \(O\) 相交于 \(B\)，\(C\)，\(D\) 为 \(BC\) 中点，若 \(PO=2\)，则 \(PA\cdot PD\) 的最大值为（\quad）

\choices
    {1}
    {\(2\)}
    {\(1+\sqrt2\)}
    {\(2+\sqrt2\)}
\end{problem}

\section{填空题}

\begin{problem}
已知点 \(A(1,5)\) 在抛物线 \(C:y^2=2px\) 上，则 \(A\) 到 \(C\) 的准线距离为\fillinblank{}.
\end{problem}

\begin{problem}
若 \(x\)，\(y\) 满足 \(x+2y\le 9\)，\(3x+y\ge 7\)，则 \(z=2x-y\) 的最大值为\fillinblank{}.
\end{problem}

\begin{problem}
已知等比数列 \(\{a_n\}\) 满足 \(a_2a_4a_5=a_3a_6\)，\(a_9a_{10}=-8\)，则 \(a_7=\)\fillinblank{}.\fillinblank{}.
\end{problem}

\begin{problem}
设 \(a\in(0,1)\)，若函数
\(f(x)=a^x+(1+a)^{x+1}\)
在 \((0,+\infty)\) 上单调递增，则 \(a\) 的取值范围是（\quad）
\end{problem}

\section{解答题}

\begin{problem}
某厂为比较甲乙两种工艺对橡胶产品伸缩率的处理效应，进行 10 次配对试验，每次配对试验选用材质相同的两个橡胶产品，随机地选其中一个用甲工艺处理，另一个用乙工艺处理，测量处理后的橡胶产品的伸缩率. 甲，乙两种工艺处理后的橡胶产品的伸缩率分别记为 \(x_i\)，\(y_i\)（\(i=1\)，\(2\)，\(\cdots\)，\(10\)）. 试验结果如下：

\begin{center}
\begin{tabular}{c|cccccccccc}
i&1&2&3&4&5&6&7&8&9&10\\\hline
\(x_i\)&545&533&551&522&575&544&541&568&596&548\\
\(y_i\)&536&527&543&530&560&533&522&550&576&536
\end{tabular}
\end{center}

\(z_i=x_i-y_i\)（\(i=1\)，\(2\)，\(\cdots\)，\(10\)），记 \(z_1\)，\(z_2\)，\(\cdots\)，\(z_{10}\) 的样本平均数为 \(\bar z\)，样本方差为 \(s^2\).

\begin{enumerate}
\item 求 \(\bar z\)，\(s^2\)；
\item 判断甲工艺处理后的橡胶产品的伸缩率较乙工艺处理后的橡胶产品的伸缩率是否有显著提高（如果 \(\bar z\ge 2\sqrt{\frac{s^2}{10}}\)，则认为甲工艺处理后的橡胶产品的伸缩率较乙工艺处理后的橡胶产品的伸缩率有显著提高，否则不认为有显著提高）.
\end{enumerate}
\end{problem}

\begin{problem}
在 \(\triangle ABC\) 中，已知 \(\angle BAC=120^\circ\)，\(AB=2\)，\(AC=1\).

\begin{enumerate}
\item 求 \(\sin \angle ABC\)；
\item 若 \(D\) 为 \(BC\) 上一点，且 \(\angle BAD=90^\circ\)，求 \(\triangle ADC\) 的面积.
\end{enumerate}
\end{problem}

\begin{problem}
如图，在三棱锥 \(P-ABC\) 中，\(AB\perp BC\)，\(AB=2\)，\(BC=2\sqrt2\)，\(PB=PC=\sqrt6\)，\(BP\)，\(AP\)，\(BC\) 的中点分别为 \(D\)，\(E\)，\(O\)，\(AD=\sqrt5 DO\)，点 \(F\) 在 \(AC\) 上，\(BF\perp AO\).



\begin{enumerate}
\item 证明：\(EF // \text{平面 }ADO\)；
\item 证明：\(\text{平面 }ADO\perp\text{平面 }BEF\)；
\item 求二面角 \(D-AO-C\) 的正弦值.
\end{enumerate}

\bitmapfigure[max width=6.0cm]{img/2023/national_paper_b_science/q19_fig1.png}
\end{problem}

\begin{problem}
已知椭圆
\(\frac{y^2}{a^2}+\frac{x^2}{b^2}=1\)（\(a>b>0\)），其离心率是 \(\frac{\sqrt5}{3}\)，点 \(A(-2,0)\) 在 \(C\) 上.

\begin{enumerate}
\item 求 \(C\) 的方程；
\item 过点 \((-2,3)\) 的直线交 \(C\) 于 \(P\)，\(Q\) 两点，直线 \(AP\)，\(AQ\) 与 \(y\) 轴的交点分别为 \(M\)，\(N\)，证明：线段 \(MN\) 的中点为定点.
\end{enumerate}
\end{problem}

\begin{problem}
已知
\(f(x)=\left( \frac1x+a \right)\ln(1+x)\).

\begin{enumerate}
\item 当 \(a=-1\) 时，求曲线 \(y=f(x)\) 在点 \((1,f(1))\) 处的切线方程；
\item 是否存在 \(a\)，\(b\)，使得曲线 \(y=f\left( \frac1x \right)\) 关于直线 \(x=b\) 对称，若存在，求 \(a\)，\(b\) 的值，若不存在，说明理由；
\item 若 \(f(x)\) 在 \((0,+\infty)\) 存在极值，求 \(a\) 的取值范围.
\end{enumerate}
\end{problem}

\section{选考题：解答题}

\begin{problem}
在直角坐标系 \(xOy\) 中，以坐标原点 \(O\) 为极点，\(x\) 轴正半轴为极轴建立极坐标系，曲线 \(C_1\) 的极坐标方程为 \(\rho=2\sin\theta\)（\(\frac{\pi}{4}\le\theta\le\frac{\pi}{2}\)），曲线 \(C_2\) 的参数方程为 \(x=2\cos\alpha\)，\(y=2\sin\alpha\)（\(\alpha\) 为参数，\(\frac{\pi}{2}<\alpha<\pi\)）.

\begin{enumerate}
\item 写出 \(C_1\) 的直角坐标方程；
\item 若直线 \(y=x+m\) 既与 \(C_1\) 没有公共点，也与 \(C_2\) 没有公共点，求 \(m\) 的取值范围.
\end{enumerate}
\end{problem}

\begin{problem}
已知 \(f(x)=2|x|+|x-2|\).

\begin{enumerate}
\item 求不等式 \(f(x)\le6-x\) 的解集；
\item 在直角坐标系 \(xOy\) 中，求由不等式 \(f(x)\le y\) 与 \(x+y-6\le0\) 所确定的平面区域的面积.
\end{enumerate}
\end{problem}

